% Created 2026-06-27 Sat 00:05
% Intended LaTeX compiler: pdflatex
\documentclass{article}
\usepackage[utf8]{inputenc}
\usepackage[T1]{fontenc}
\usepackage{graphicx}
\usepackage{longtable}
\usepackage{wrapfig}
\usepackage{rotating}
\usepackage[normalem]{ulem}
\usepackage{amsmath}
\usepackage{amssymb}
\usepackage{capt-of}
\usepackage{hyperref}

\usepackage{fancyhdr}
\usepackage[top=25mm, left=25mm, right=25mm]{geometry}
\usepackage{longtable}
\fancyhead[R]{}
\setlength\headheight{43.0pt}
\author{Emilio Albán, Ariel Montufar, Samuel Churaco, Karen Suarez}
\date{2026-06-19}
\title{Taller de Comandos Linux}
\hypersetup{
 pdfauthor={Emilio Albán, Ariel Montufar, Samuel Churaco, Karen Suarez},
 pdftitle={Taller de Comandos Linux},
 pdfkeywords={},
 pdfsubject={},
 pdfcreator={Emacs 27.1 (Org mode 9.7.5)}, 
 pdflang={Español}}
\begin{document}

\maketitle
\fancyhead[C]{\includegraphics[scale=0.05]{../images/logoEPN.jpg}\\
ESCUELA POLITÉCNICA NACIONAL\\FACULTAD DE INGENIERÍA DE SISTEMAS\\
ARQUITECTURA DE COMPUTADORES}
\thispagestyle{fancy}


\section{Objetivos}
\label{sec:orgf409a0e}

\begin{itemize}
\item Conocer el funcionamiento del terminal de \emph{Linux}
\item Aprender a manejar comandos del terminal de \emph{Linux}
\end{itemize}

\section{Instrucciones}
\label{sec:org0f5000a}
Este taller/tutorial de comandos de Linux tiene por objetivo ejercitar
en la aplicación de comandos estudiados en clase. Para esto siga las
Instrucciones indicadas en cada uno de los numerales. Ejecute los
ejercicios llenando los comandos adecuados dentro de los bloques de
código de Emacs Org. Los resultados de la ejecución se obtienen
haciendo \texttt{C-c C-c}

Para todos los ejercicios se sugiere usar directamente un
distro de \emph{Linux} o el \emph{WSL} configurado para la clase de
\emph{Arquitectura de Computadores}. Ejecute los comandos en la terminal de
Linux. Copie los comandos usados a este archivo.

Si se ha utilizado IA para consultar temas referentes al desarrollo de
código en el punto \ref{sec:orgcc5575e}, adjunte el archivo \texttt{.md} con el
reporte de los diferentes prompts y respuestas dadas por la IA.

\subsection{Ejecución desde Emacs}
\label{sec:org7b45ed5}

Si desea ejecutar los comandos desde Emacs debe observar que los
comandos del \texttt{shell} de Linux se ejecutarán desde la ubicación del
archivo \texttt{.org}. Por tanto, controle la ubicación de \texttt{home} en los
bloques de código. Para ejecutar los bloques de código ejecute: \texttt{C-c C-c}. 

\section{Comandos de Linux}
\label{sec:orgfe75a79}
Para este taller descargaremos los datos del \emph{Fashion-Mnist}, un
dataset popular para la prueba de algoritmos de machine learning. Al
final, el dataset quedará organizado para ser aprovechado para
entrenar un modelo.

\begin{enumerate}
\item Utilice \texttt{cd} para localizarse en la posición de home de su usuario y
verifique ejecutando el comando \texttt{pwd}

\begin{verbatim}
   cd
   pwd
\end{verbatim}

\begin{verbatim}
/home/arielmontufar
\end{verbatim}
\end{enumerate}



\begin{enumerate}
\item Cree tres directorios denominados FashionData, Train y
Test. Verifique usando el comando \texttt{ls}
\begin{verbatim}
   ls -ld ~/FashionData ~/Train ~/Test
\end{verbatim}

\begin{verbatim}
drwxr-xr-x 2 arielmontufar arielmontufar 4096 Jun 25 21:59 /home/arielmontufar/FashionData
drwxr-xr-x 2 arielmontufar arielmontufar 4096 Jun 25 22:13 /home/arielmontufar/Test
drwxr-xr-x 2 arielmontufar arielmontufar 4096 Jun 25 21:59 /home/arielmontufar/Train
\end{verbatim}
\end{enumerate}


\begin{enumerate}
\item Investigue sobre el comando \texttt{wget} y utilícelo para descargar en
\textasciitilde{}/FashionData los datos de train y test del /Fashion-Mnist desde el
siguiente enlace: \href{https://github.com/zalandoresearch/fashion-mnist}{fashion>-mnist}
\begin{enumerate}
\item Imágenes de entrenamiento:

\url{http://fashion-mnist.s3-website.eu-central-1.amazonaws.com/train-images-idx3-ubyte.gz}
\item Etiquetas de entrenamiento:

\url{http://fashion-mnist.s3-website.eu-central-1.amazonaws.com/train-labels-idx1-ubyte.gz}
\item Imágenes de prueba:

\url{http://fashion-mnist.s3-website.eu-central-1.amazonaws.com/t10k-images-idx3-ubyte.gz}
\item Etiquetas de prueba:

\url{http://fashion-mnist.s3-website.eu-central-1.amazonaws.com/t10k-labels-idx1-ubyte.gz}
\end{enumerate}
\begin{verbatim}

cd ~/FashionData
wget -c http://fashion-mnist.s3-website.eu-central-1.amazonaws.com/train-images-idx3-ubyte.gz
wget -c http://fashion-mnist.s3-website.eu-central-1.amazonaws.com/train-labels-idx1-ubyte.gz
wget -c http://fashion-mnist.s3-website.eu-central-1.amazonaws.com/t10k-images-idx3-ubyte.gz
wget -c http://fashion-mnist.s3-website.eu-central-1.amazonaws.com/t10k-labels-idx1-ubyte.gz

ls -lh
\end{verbatim}

\begin{verbatim}
total 30M
-rw-r--r-- 1 arielmontufar arielmontufar 4.3M Aug 31  2017 t10k-images-idx3-ubyte.gz
-rw-r--r-- 1 arielmontufar arielmontufar 5.1K Aug 31  2017 t10k-labels-idx1-ubyte.gz
-rw-r--r-- 1 arielmontufar arielmontufar  26M Aug 31  2017 train-images-idx3-ubyte.gz
-rw-r--r-- 1 arielmontufar arielmontufar  29K Aug 31  2017 train-labels-idx1-ubyte.gz
\end{verbatim}

\item Utilice el comando \texttt{mv} de manera recursiva para pasar los datos de
entrenamiento a la carpeta \textasciitilde{}/Train
\begin{verbatim}
mv ~/FashionData/*train* ~/Train/
ls -lh ~/Train/

\end{verbatim}

\begin{verbatim}
total 51M
-rw-r--r-- 1 arielmontufar arielmontufar 26M Aug 31  2017 train-images-idx3-ubyte.gz
-rw-r--r-- 1 arielmontufar arielmontufar 26M Aug 31  2017 train-images-idx3-ubyte.gz.1
-rw-r--r-- 1 arielmontufar arielmontufar 29K Aug 31  2017 train-labels-idx1-ubyte.gz
\end{verbatim}

\item Utilice el comando \texttt{mv} de manera recursiva para pasar los datos de
prueba a la carpeta \textasciitilde{}/Test
\begin{verbatim}
mv ~/FashionData/t10k* ~/Test/
ls -lh ~/Test/
\end{verbatim}

\begin{verbatim}
total 4.3M
-rw-r--r-- 1 arielmontufar arielmontufar 4.3M Aug 31  2017 t10k-images-idx3-ubyte.gz
-rw-r--r-- 1 arielmontufar arielmontufar 5.1K Aug 31  2017 t10k-labels-idx1-ubyte.gz
\end{verbatim}

\item Verifique que la carpeta FashionData está vacía
\begin{verbatim}
ls -lh ~/FashionData/
\end{verbatim}

\begin{verbatim}
total 0
\end{verbatim}

\item Los archivos anteriores están comprimidos. Investigue como
descomprimir usando el comando \texttt{gunzip}. Si es necesario realice la
instalación del comando. Descomprima los archivos en sus
respectivas carpetas. Apunte el comando utilizado para la descompresión.
\begin{verbatim}
gunzip -d ~/Train/*.gz
gunzip -d ~/Test/*.gz
ls -lh ~/Train/
ls -lh ~/Test/
\end{verbatim}

\begin{verbatim}
total 45M
-rw-r--r-- 1 arielmontufar arielmontufar 45M Aug 31  2017 train-images-idx3-ubyte
-rw-r--r-- 1 arielmontufar arielmontufar 59K Aug 31  2017 train-labels-idx1-ubyte
total 7.5M
-rw-r--r-- 1 arielmontufar arielmontufar 7.5M Aug 31  2017 t10k-images-idx3-ubyte
-rw-r--r-- 1 arielmontufar arielmontufar 9.8K Aug 31  2017 t10k-labels-idx1-ubyte
\end{verbatim}

\item Mueva recursivamente las carpetas \textasciitilde{}/Train y \textasciitilde{}/Test dentro de \textasciitilde{}/FashionData
\begin{verbatim}
mv ~/Train ~/FashionData/
mv ~/Test ~/FashionData/
ls -R ~/FashionData/
\end{verbatim}

\begin{verbatim}
   /home/arielmontufar/FashionData/:
   Test
   Train

   /home/arielmontufar/FashionData/Test:
   t10k-images-idx3-ubyte
   t10k-labels-idx1-ubyte

   /home/arielmontufar/FashionData/Train:
   train-images-idx3-ubyte
   train-labels-idx1-ubyte
\end{verbatim}

\item Verifique la estructura de la carpeta \texttt{/FashionData usando el
   comando \textasciitilde{}tree}
\begin{verbatim}
tree ~/FashionData/
\end{verbatim}

\begin{verbatim}
/home/arielmontufar/FashionData/
├── Test
│   ├── t10k-images-idx3-ubyte
│   └── t10k-labels-idx1-ubyte
└── Train
    ├── train-images-idx3-ubyte
    └── train-labels-idx1-ubyte

3 directories, 4 files
\end{verbatim}
\end{enumerate}



\section{Problema}
\label{sec:orgcc5575e}
Se desea realizar un registro climatológico de la ciudad de
Quito. Para esto, escriba un script de Python/Java que permita obtener
datos climatológicos desde el API de
\href{https://openweathermap.org/current\#one}{openweathermap}. Considere la latitud y la longitud de Quito como
(-0.2299, -78.5249), respectivamente. Los resultados obtenidos de
la consulta al API se escriben en un archivo
\emph{clima-quito-hoy.csv}. Cada ejecución del script debe almacenar
nuevos datos en el archivo. Investigue sobre \textbf{crontab} y utilícelo
para obtener datos del API de \emph{openweathermap} cada 5 minutos
mediante la ejecución de un archivo ejecutable denominado
\emph{get-weather.sh}. Verifique los resultados. Todas las operaciones
se realizan en Linux o en el WSL. Las etapas del problema se
subdividen en:

\begin{enumerate}
\item Crear su API gratuito en \href{https://openweathermap.org/current\#one}{openweathermap}
API: f0a67d31e3a1532c0e0404b8b5b2525f
\item Escribir un script en Python/Java que realice la consulta al API y
escriba los resultados en \emph{clima-quito-hoy.csv}
\begin{enumerate}
\item Primero se creó el archivo clima-quito-hoy.csv en la terminal
con touch.
\item Desarrollo del Script:
Importante: Se utilizó la librería externa JSON, por lo tanto se
instaló el json-20231013.jar, se lo compiló y ejecutó con el
.jar ya mencionado (ya compilado y ejecutando en la terminal de
esta manera):
\end{enumerate}
\end{enumerate}

\begin{verbatim}
javac -cp .:json-20231013.jar RegistroClimaticoQuito.java
java -cp .:json-20231013.jar RegistroClimaticoQuito
\end{verbatim}

\begin{verbatim}
import java.net.URL;
import java.net.HttpURLConnection;
import java.io.*;
import org.json.*;
import java.time.LocalDateTime;

public class RegistroClimaticoQuito{

    //API KEY generado de OpenWeatherMap
    private static final String API_KEY = "f0a67d31e3a1532c0e0404b8b5b2525f";
    //Latitud y Longitud de Quito
    private static final double LATITUD = -0.2299;
    private static final double LONGITUD = -78.5249;

    //Nombre del archivo en donde se escribiran los resultados
    private static final String ARCHIVO = "clima-quito-hoy.csv";

    public static void main(String[] args) {
	try{ //Controlar errores y que el programa no quede sin finalizar
	    String url = construirURL(); //Se construye el URL Del API
	    String respuesta = consultarAPI(url); //Se consulta el API y se obtiene la respuesta JSON Como texto
	    JSONObject datos = convertirJSON(respuesta); //Se convierte de String-JSON a objeto JSON
	    guardarCSV(datos); //Se guardan los datos en clima-quito-hoy.csv

            System.out.println("Datos guardados correctamente.");

        }catch (Exception e) {
            System.out.println("Error: " + e.getMessage());
        }
    }

    public static String construirURL(){
        // se arma la URL con latitud, longitud de Quito, unidades(C°) y API key
        String url =
                "https://api.openweathermap.org/data/2.5/weather"
                + "?lat=" + LATITUD
                + "&lon=" + LONGITUD
                + "&units=metric"
                + "&lang=es"
                + "&appid=" + API_KEY;

        // se retorna la URL completa
        return url;
    }

    public static String consultarAPI(String urlString) throws Exception {

        URL url = new URL(urlString);// crear objeto URL apartir del string URL
        HttpURLConnection con = (HttpURLConnection) url.openConnection(); // abrir conexión HTTP
        con.setRequestMethod("GET");// definir método GET
	
        // Se lee respuesta del servidor
        BufferedReader reader = new BufferedReader(new InputStreamReader(con.getInputStream()));

        StringBuilder respuesta = new StringBuilder();// almacenar respuesta completa
        String linea; // variable para lectura línea por línea
        // leer todo el contenido de la respuesta
        while ((linea = reader.readLine()) != null) {
            respuesta.append(linea);
        }
        
        reader.close();// cerrar lector
        return respuesta.toString();// devolver JSON como texto
    }

    public static JSONObject convertirJSON(String respuesta) {
        // convertir string JSON a objeto JSON para los datos al .csv
        JSONObject obj = new JSONObject(respuesta);
        // retornar objeto JSON
        return obj;
    }

    public static void guardarCSV(JSONObject obj) throws Exception {

        // verifica si el archivo ya existe cuando se coloque el encabezado
        File file = new File(ARCHIVO);
        FileWriter csvWriter = new FileWriter(ARCHIVO, true); //se abre el archivo en modo append (no borra datos anteriores)

        // si el archivo está vacío/nuevo,se escribe el encabezado
        if (file.length() == 0) {
            csvWriter.append("fecha,ciudad,temperatura,humedad,presion,descripcion\n");
        }

	String fechaHora = LocalDateTime.now().toString(); //Fecha de guardado
        //SE EXTRAE DATOS DEL JSON:
        String ciudad = obj.getString("name"); //el nombre de la ciudad
        JSONObject main = obj.getJSONObject("main"); // el conjunto main del JSON de lo retornado a la llamada API (dado que contiene temp, feels_like, temp_min, etc..)
        double temp = main.getDouble("temp"); // temperatura del conjunto main
        int humedad = main.getInt("humidity");// humedad del conjunto main
        int presion = main.getInt("pressure");// presión del conjunto main
        JSONArray weather = obj.getJSONArray("weather");//el arreglo weather del JSON de lo retornado a la llamada API
        String descripcion = weather.getJSONObject(0).getString("description");//descripción del clima (guardado en el primer valor, por eso el 0)
	
        // ESCRIBIR EN CSV usando coma como separador
        csvWriter.append(fechaHora + "," +
                         ciudad + "," +
                         temp + "," +
                         humedad + "," +
                         presion + "," +
                         descripcion + "\n");
        csvWriter.close();//se cierra el archivo
    }
}
\end{verbatim}

\begin{enumerate}
\item Desarrollar un ejecutable \emph{get-weather.sh} para ejecutar el
programa Python/Java.
Nota: Se creó el .sh en la terminal con touch y su
contenido es el siguiente:

\begin{verbatim}
   #!/bin/bash
# 1. Ubica el script en el directorio del proyecto
cd "$(dirname "$0")"

# 2. Se compila el archivo Java
javac -cp .:json-20231013.jar RegistroClimaticoQuito.java

# 3. Ejecuta el programa Java usando el classpath especificado
java -cp .:json-20231013.jar RegistroClimaticoQuito

# 4. Dar permisos de ejecución pero directamente desde la terminal con: chmod +x get-weather.sh
\end{verbatim}
\end{enumerate}


\begin{enumerate}
\item Configurar Crontab para la adquisición de datos. Escriba el comando
configurado.

\begin{verbatim}
crontab -e
# Se agrega directamente a la linea del crontab:
# */5 * * * * $HOME/G1ArqComp/G1ArqComp/Talleres/get-weather.sh
# significa cada 5 minutos el cron ejecutará el programa.
\end{verbatim}
\end{enumerate}
\end{document}
